\documentclass[12pt]{article}
\usepackage[utf8]{inputenc}
\usepackage[T1]{fontenc}
\usepackage[margin=1in]{geometry}
\usepackage{setspace}
\usepackage{xcolor}
\usepackage[colorlinks=true, urlcolor=blue!60!black, linkcolor=blue!60!black]{hyperref}
\setlength{\parindent}{2em}
\setlength{\parskip}{1.25em}
\renewcommand{\baselinestretch}{1.0}
\begin{document}
\begin{center}
{\LARGE \textit{Tandon v. Newsom}, 593 U.S. 61 (2021)\par}
\vspace{0.8em}
{\large Case Brief --- Four-Page Edition\par}
\end{center}
\vspace{1em}

\section{Memory Jogger}

In a 5--4 emergency-docket \textit{per curiam} order, the Court enjoined California's three-household cap on in-home religious gatherings, adopting the ``most-favored-nation'' rule: treating \textit{any} comparable secular activity more favorably than religious exercise triggers strict scrutiny under the Free Exercise Clause (\href{https://www.supremecourt.gov/opinions/20pdf/20a151_4g15.pdf}{Slip Opinion}).

\section{Detailed Case Facts}

During the COVID-19 pandemic, California's ``Blueprint for a Safer Economy'' capped both religious and secular gatherings in private homes at three households. At the same time, it let hair salons, retail stores, personal care services, movie theaters, private suites at sporting events and concerts, and indoor restaurants bring together more than three households under sector-specific protocols. In October 2020, California residents, including Ritesh Tandon and religious plaintiffs Pastor Jeremy Wong and Karen Busch, sued. Wong and Busch had regularly hosted Bible studies and collective prayer for eight to twelve people in their homes. After the emergency application was filed, California announced changes effective April 15, 2021, while keeping the restrictions in force until then.

\section{Procedural History}

Judge Lucy H. Koh (N.D. Cal.) denied a preliminary injunction, relying on uncontested expert testimony that in-home gatherings posed greater risks than commercial venues and finding that the cap applied equally to all in-home gatherings. A divided Ninth Circuit panel (Judges Milan D. Smith, Jr., and Bridget S. Bade; Judge Patrick J. Bumatay dissenting in part) denied an injunction pending appeal, comparing religious home gatherings only with secular home gatherings and applying rational-basis review and noting that commercial precautions might not ``translate readily'' to private homes. On April 9, 2021, without argument or merits briefing, the Supreme Court granted an injunction pending appeal. The district court later entered a stipulated permanent injunction with \$500,000 in attorney's fees.

\section{Judicial Votes}

Application granted, \textbf{5--4}.
\begin{itemize}
\item \textbf{Per curiam majority:} Justices Thomas, Alito, Gorsuch, Kavanaugh, and Barrett.
\item \textbf{Dissent without opinion:} Chief Justice Roberts would have denied the application.
\item \textbf{Dissent:} Justice Kagan, joined by Justices Breyer and Sotomayor.
\end{itemize}

\section{Holding}

California was enjoined from enforcing the cap against the applicants' in-home religious exercise. Under \textit{Employment Division v. Smith}: (1) regulations are not neutral and generally applicable, and trigger strict scrutiny, whenever they treat \textit{any} comparable secular activity more favorably than religious exercise, even if other comparable secular activities are treated equally or worse; (2) comparability is judged against the asserted government interest, not the reasons people gather; (3) the government must show that precautions allowed for secular activities would not suffice for religious exercise; and (4) changing a restriction during litigation does not moot the case when reinstatement remains a threat. The admonition that this was the fifth summary reversal of the Ninth Circuit is dicta.

\section{Analysis of Opinions}

\textbf{Per curiam.} The opinion consolidated \textit{Roman Catholic Diocese of Brooklyn v. Cuomo} and \textit{South Bay United Pentecostal Church v. Newsom} into a four-point framework. Treating secular home gatherings the same as religious ones was ``no answer.'' The Ninth Circuit erred in rejecting commercial comparators simply because earlier cases involved public buildings, and in speculating that commercial precautions might not translate to homes. Nor was the case moot: the cap remained in force until April 15, and officials with a track record of changing the rules retained authority to reinstate heightened restrictions. The applicants were likely to succeed, faced irreparable harm from even brief losses of free exercise, and the State had not shown that less restrictive precautions would imperil public health.

\textbf{Roberts (dissenting without opinion).} He did not join Justice Kagan's dissent. Consistent with his \textit{South Bay} concurrences, he had favored deference to accountable officials in a public-health emergency unless a State singled out houses of worship.

\textbf{Kagan (dissenting).} The obvious comparator for at-home religious gatherings is at-home secular gatherings, which California capped identically. Even on the majority's frame, the record showed three risk differences: home gatherings involve longer interactions, homes are smaller and less ventilated, and masking and distancing are harder to enforce. Ignoring those undisputed findings impaired the State's scientific response to an emergency.

\section{Examples: Future Applications}

\begin{enumerate}
\item \textbf{Same-Side: Police Beards.} A department bans beards but grants medical exemptions while denying Muslim officers religious ones. The medical exception undermines uniformity equally, triggering strict scrutiny.
\item \textbf{Same-Side: Religious Student Club.} A university applies an ``all-comers'' rule to a religious club while letting secular affinity clubs restrict membership. Exempting comparable secular groups triggers strict scrutiny against denying the religious club an exemption.
\item \textbf{Opposite-Side: Noise Ordinance.} A residential noise limit with no secular exemptions does not trigger strict scrutiny against amplified midnight prayers, because no comparable secular activity receives more favorable treatment.
\item \textbf{Opposite-Side: Venomous Snakes.} A ban on private possession of venomous snakes exempts accredited zoos. If zoos pose materially lower risk, they are not comparable to private snake handling in worship.
\item \textbf{Fence-Sitter: Vaccine Mandate.} A hospital-worker mandate allows medical but not religious exemptions. Lower courts divide over whether the two pose identical risks or whether medical exemptions are fewer, verifiable, and serve worker health.
\end{enumerate}

\section{Critique}

\textbf{Praise and methodological doubts.} Called ``the most important free exercise decision since 1990,'' \textit{Tandon} adopted Douglas Laycock's most-favored-nation reading without overruling \textit{Smith}. Supporters say it stops governments from subordinating religion to favored commerce; methodological critics note that it reshaped doctrine in a four-page emergency order with no historical analysis.

\textbf{Progressive and procedural critique.} Jim Oleske, Alan Brownstein, Vikram Amar, and Nelson Tebbe argue that the rule hollows out \textit{Smith}, since nearly every statute has secular exceptions, and leaves unresolved the level of generality and burden on comparability. Steve Vladeck faults the Court for a landmark shift on the ``shadow docket'' without briefing, argument, or deference to fact-finding.

\textbf{Competing values.} The case pits the commitment that religious conscience may not be quietly devalued below secular interests during crises against institutional competence, judicial restraint, and States' capacity to design risk-calibrated rules without triggering strict scrutiny whenever a law accommodates a single secular necessity.

\section{Key Quotations}

\begin{quote}
``First, government regulations are not neutral and generally applicable, and therefore trigger strict scrutiny under the Free Exercise Clause, whenever they treat \textit{any} comparable secular activity more favorably than religious exercise.''
\end{quote}
\noindent --- \textit{Per Curiam}, 593 U.S. at 62.

\begin{quote}
``Comparability is concerned with the risks various activities pose, not the reasons why people gather.''
\end{quote}
\noindent --- \textit{Per Curiam}, 593 U.S. at 62.

\begin{quote}
``As the per curiam's reliance on separate opinions and unreasoned orders signals, the law does not require that the State equally treat apples and watermelons.''
\end{quote}
\noindent --- Justice Kagan, dissenting, 593 U.S. at 65.

\end{document}
